%%%%%%%%%%  scan idx 470  |  folio 457  %%%%%%%%%%
\SGAPrintPage{457}%
The ideal generated by \(P\) is a \(\lambda\)-ideal. To see this, first observe that the relations \(\lambda^i(N)=0\) for \(i>d\) imply that the \(\lambda^j(N-L)\) for \(j\ge d\) are multiples of \(\lambda^d(N-L)\). Working in \(1+\widehat{K[L]/(P)[[t]]}^+\), the image of \(\lambda_t(N-L)\) is a polynomial of degree \(<d\), and a universal argument then shows that its \(\lambda^i\) for \(i\ge d\) equal \(1\) (cf. V 4.8). Since in \(1+\widehat{K[L][[t]]}^+\) we have \(\lambda_t(\lambda^d(N-L))=\lambda^d(\lambda_t(N-L))\), we deduce that the \(\lambda^i(\lambda^d(N-L))\) for \(i\ge1\) are multiples of \(P\). By the relations V 2.4.2, this implies that the ideal generated by \(P\) is a \(\lambda\)-ideal.

Put then \(K_1=K[L]/(P)\). We denote the natural homomorphism \(K\to K_1\) by \(g^*\); it is a \(\lambda\)-homomorphism. Furthermore, \(K_1\) is a free \(K\)-module with basis the elements \(1,L,\ldots,L^{d-1}\). The relation \(\lambda^d(N)=\lambda^{d-1}(N-L).L\) and the invertibility hypothesis on \(\lambda^d(N)\) imply that \(L\) is invertible. Consequently, \(K_1\) has as a basis over \(K\) the elements \(1,L^{-1},\ldots,L^{-d+1}\); we then define a map \(g_*:K_1\to K\) by assigning to every element \(y'=y_0+y_1L^{-1}+\ldots+y_{d-1}L^{-d+1}\) of \(K_1\) the element \(y_0\) of \(K\). Clearly the map thus defined is an additive, \(K\)-linear homomorphism, satisfying \(g_*\circ g^*=\Id\).

Next we endow \(K_1\times K'\) with a \(\lambda\)-ring structure, putting:
\sgaStableTarget{sga6.e07.eq.4-7-1}\sgaStableTarget{sga6.e07.eq.4-7-2}\begin{align}
 (y'_1,x_1)*(y'_2,x_2)
   &=(y'_1y'_2(1-L)+x_1y'_2+x_2y'_1,\; x_1x_2); \tag{4.7.1}\\
 \lambda^n_*(y',x)
   &=\left(\sum_{p=1}^{n}\lambda^{\,n-p}(x)\lambda^p(L,y'),\;
     \lambda^n(x)\right). \tag{4.7.2}
\end{align}
In these formulas, \(K_1\) is regarded as a \(K'\)-algebra through the homomorphism \(g^*\circ i^*\).
Next, for \(y_1,y_2\in K\) and \(n\ge1\), put:
\sgaStableTarget{sga6.e07.eq.4-7-3}\sgaStableTarget{sga6.e07.eq.4-7-4}\begin{align}
 y_1*y_2&=-\,y_1y_2\lambda_{-1}(N); \tag{4.7.3}\\
 \lambda^n_*(y_1)&=-\,\lambda^n(N,-y_1); \tag{4.7.4}
\end{align}

%%%%%%%%%%  scan idx 471  |  folio 458  %%%%%%%%%%
\SGAPrintPage{458}%
and define a homomorphism:
\[
  u:K\longrightarrow K_1\times K'
\]
by putting:
\[
  u(y)=\bigl(g^*(y).\lambda_{-1}(F),-i_*(y)\bigr),
\]
with \(F=N-L\).

\sgaStableTarget{sga6.e07.prop.4-7-5}\textbf{Proposition 4.7.5.} \uline{The operations defined by} (4.7.1) \uline{and} (4.7.2) \uline{endow} \(K_1\times K'\) \uline{with a} \(\lambda\)-\uline{ring structure}, \uline{and} \(u\) \uline{is an injective homomorphism identifying} \(K\), \uline{equipped with the operations defined by} (4.7.3) \uline{and} (4.7.4), \uline{with a} \(\lambda\)-\uline{ideal of} \(K_1\times K'\).

It is immediate to verify that \(K_1\times K'\) is a \(\lambda\)-ring: one can refer, for example, to V 5.5, where an entirely analogous construction is given. The second part of the proposition is also proved without difficulty using the preceding relations, the relation \(i_*(y_1y_2\lambda_{-1}(N))=i_*(y_1).i_*(y_2)\) which follows immediately from the hypotheses, the relation \(g_*(\lambda_{-1}(F))=1\) (see VI 5.9), and the relations of V 5.

Put then:
\[
  K'_1=(K_1\times K')/K
\]
and denote by \(v\) the canonical homomorphism \(K_1\times K'\to K'_1\). Regard \(K'_1\) as a \(\lambda\)-ring through the quotient structure.

Next define additive homomorphisms:
\[
\begin{array}{llll}
j_*:K_1\to K'_1; & j^*:K'_1\to K_1; \\
f_*:K'_1\to K'; & f^*:K'\to K'_1 .
\end{array}
\]

The homomorphism \(j_*\) is the composite \(K_1\to K_1\times K'\to K'_1\). Likewise, the homomorphism \(f^*\) is the composite \(K'\to K_1\times K'\to K'_1\). We then have:
\[
  v=j_*+f^*.
\]

%%%%%%%%%%  scan idx 472  |  folio 459  %%%%%%%%%%
\SGAPrintPage{459}%
Next define a homomorphism \(K_1\times K'\to K_1\) by assigning to \((y',x)\) the element \(y'(1-L)+g^*\circ i^*(x)\); this homomorphism vanishes on \(K\), and consequently defines a homomorphism \(j^*:K'_1\longrightarrow K_1\). Finally, define a homomorphism \(K_1\times K'\to K'\) by assigning to \((y',x)\) the element \(i_*\circ g_*(y')+x\); this homomorphism vanishes on \(K\) and therefore defines a homomorphism \(f_*:K'_1\longrightarrow K'\).

\sgaStableTarget{sga6.e07.prop.4-7-6}\textbf{Proposition 4.7.6.}
\begin{enumerate}[label=\roman*)]
\item \(j^*\) \uline{and} \(f^*\) \uline{are} \(\lambda\)-\uline{homomorphisms};
\item \(j_*\) \uline{is} \(K'_1\)-\uline{linear and} \(f_*\) \uline{is} \(K'\)-\uline{linear};
\item \(j^*\circ j_*\) \uline{is the homomorphism} ``\uline{multiplication by} \(1-L\)'';
\item \uline{for every} \(y'\in K_1\), \uline{and every} \(p\ge1\), \(j_*\bigl(\lambda^p(L,y')\bigr)=\lambda^p(j_*(y'))\);
\item \(f_*\circ f^*=\Id\).
\end{enumerate}

We need only set the machinery in motion.

Returning to the geometric situation recalled above, we see that \(K^{\bullet}(X')\) is constructed from \(K^{\bullet}(Y)\), \(K^{\bullet}(X)\), the maps \(i^*\) and \(i_*\), and the element \(N\) of \(K^{\bullet}(Y)\) by the preceding method; the \(\lambda\)-operations on \(K^{\bullet}(X')\), however, agree with the \(\lambda\)-operations defined above on \(K'_1\) only up to torsion, as long as relation (4.3.1) is proved only modulo torsion.

\sgaStableTarget{sga6.e07.thm.4-8}\textbf{Theorem 4.8.} \uline{Let} \(i:Y\to X\) \uline{be a regular closed immersion of constant codimension} \(d\), \(X'\) \uline{the scheme obtained by blowing up} \(Y\), \uline{and} \(Y'=X'\times_XY\), \uline{so that we have the commutative diagram}:
\[
\begin{tikzcd}[column sep=large,row sep=large]
Y' \arrow[r,"j"] \arrow[d,"g"'] & X' \arrow[d,"f"]\\
Y \arrow[r,"i"] & X.
\end{tikzcd}
\]
\uline{Suppose that} \(X\) \uline{is quasi-compact and has an ample sheaf}. \uline{Then for every} \(k\) \uline{we have the exact sequence}:

%%%%%%%%%%  scan idx 473  |  folio 460  %%%%%%%%%%
\SGAPrintPage{460}%
\[
0\longrightarrow
\operatorname{Gr}^{k-d}K^{\bullet}(Y)_{\mathbb Q}
\xrightarrow{u^{\gr}}
\operatorname{Gr}^{k-1}K^{\bullet}(Y')_{\mathbb Q}
\times
\operatorname{Gr}^{k}K^{\bullet}(X)_{\mathbb Q}
\xrightarrow{v^{\gr}}
\operatorname{Gr}^{k}K^{\bullet}(X')_{\mathbb Q}
\]
\uline{where we put}, \uline{for} \(y\in \operatorname{Gr}^{k-d}K^{\bullet}(Y)_{\mathbb Q}\):
\[
  u^{\gr}(y)=
  \bigl(g^{\gr}(y).(-1)^{d-1}c^{d-1}(F),-i_{\gr}(y)\bigr),
\]
\uline{and for} \(y'\in \operatorname{Gr}^{k-1}K^{\bullet}(Y')_{\mathbb Q}\) \uline{and} \(x\in \operatorname{Gr}^{k}K^{\bullet}(X)_{\mathbb Q}\):
\[
  v^{\gr}(y',x)=j_{\gr}(y')+f^{\gr}(x),
\]
\uline{with the notation of} 3.1 \uline{and} 3.3, \uline{and} \(c^{d-1}(F)\) \uline{denoting the} \((d-1)\)-\uline{st Chern class of} \(F\) (V 6.7). \uline{Moreover}, \(u^{\gr}\) \uline{has a left inverse} \(u'^{\gr}\), \uline{defined by} \(u'^{\gr}(y',x)=g_{\gr}(y')\).

\uline{If in addition} \(X\) \uline{is noetherian and} \(Y\) \uline{regular}, \uline{we have the exact sequence}:
\[
0\longrightarrow
\operatorname{Gr}^{k-d}K^{\bullet}(Y)_{\mathbb Q}
\xrightarrow{u^{\gr}}
\operatorname{Gr}^{k-1}K^{\bullet}(Y')_{\mathbb Q}
\times
\operatorname{Gr}^{k}K^{\bullet}(X)_{\mathbb Q}
\xrightarrow{v^{\gr}}
\operatorname{Gr}^{k}K^{\bullet}(X')_{\mathbb Q}
\longrightarrow0
\]

\emph{Proof.} Observe that, when \(Y\) is regular, one could give a purely algebraic proof in terms of the theory of \(\lambda\)-rings, using the description of the situation in 4.7. We shall instead follow the proof of 3.7.

We know that the ``inverse-image'' homomorphisms \(f^*,g^*\), being \(\lambda\)-homomorphisms, are compatible with the \(\lambda\)-ring filtrations, and consequently define homomorphisms on the associated gradeds; hence the homomorphisms:
\[
  f^{\gr}:\operatorname{Gr}^kK^{\bullet}(X)\to\operatorname{Gr}^kK^{\bullet}(X');
  \qquad
  g^{\gr}:\operatorname{Gr}^{k-d}K^{\bullet}(Y)\to\operatorname{Gr}^{k-d}K^{\bullet}(Y'),
\]
and the analogous homomorphisms obtained after tensoring with \(\mathbb Q\). Furthermore, 4.6 shows that after tensoring with \(\mathbb Q\), the homomorphisms \(i_*,j_*\) pass to the associated gradeds, and therefore define:
\[
  i_{\gr}:\operatorname{Gr}^{k-d}K^{\bullet}(Y)_{\mathbb Q}
  \to \operatorname{Gr}^{k}K^{\bullet}(X)_{\mathbb Q};
  \qquad
  j_{\gr}:\operatorname{Gr}^{k-1}K^{\bullet}(Y')_{\mathbb Q}
  \to \operatorname{Gr}^{k}K^{\bullet}(X')_{\mathbb Q}.
\]

%%%%%%%%%%  scan idx 474  |  folio 461  %%%%%%%%%%
\SGAPrintPage{461}%
Finally, observe that \(\lambda_{-1}(F)=(-1)^{d-1}\gamma^{d-1}(F)\), so the morphisms \(u^{\gr}\) and \(v^{\gr}\) are obtained by passage to the quotient from the morphisms \(u\) and \(v\) of 3.7. Thus we have \(v^{\gr}\circ u^{\gr}=0\).

By VI 5.8, \(g_*(\Fil^{k-1}K^{\bullet}(Y'))\subset \Fil^{k-d}K^{\bullet}(Y)\) for every \(k\), since \(Y'=\mathbf P(\underline{N})\), with \(\underline{N}\) of rank \(d\). We therefore obtain a homomorphism:
\[
  g_{\gr}:\operatorname{Gr}^{k-1}K^{\bullet}(Y')
  \to\operatorname{Gr}^{k-d}K^{\bullet}(Y),
\]
and an analogous homomorphism after tensoring with \(\mathbb Q\). The homomorphism \(u'^{\gr}\) is therefore obtained from \(u'\) (3.7) by passage to the associated graded. Since \(u'\) is a left inverse of \(u\), \(u'^{\gr}\) is a left inverse of \(u^{\gr}\), proving at the same time the injectivity of \(u^{\gr}\).

To show that \(\Ker(v^{\gr})=\Im(u^{\gr})\), we must show that for \(y'\in\Fil^{k-1}K^{\bullet}(Y')_{\mathbb Q}\), and \(x\in\Fil^kK^{\bullet}(X)_{\mathbb Q}\), satisfying:
\sgaStableTarget{sga6.e07.eq.4-6-1}\begin{equation}\tag{4.8.1}
  j_*(y')+f^*(x)\in\Fil^{k+1}K^{\bullet}(X')_{\mathbb Q},
\end{equation}
there exists \(y\in\Fil^{k-d}K^{\bullet}(Y)_{\mathbb Q}\) such that:
\[
\begin{aligned}
\text{i)}\quad &x+i_*(y)\in\Fil^{k+1}K^{\bullet}(X)_{\mathbb Q},\\
\text{ii)}\quad &y'-g^*(y)\lambda_{-1}(F)\in\Fil^{k}K^{\bullet}(Y')_{\mathbb Q}.
\end{aligned}
\]

Put:
\[
  y=g_*(y').
\]
Then by VI 5.8, \(y\in\Fil^{k-d}K^{\bullet}(Y)_{\mathbb Q}\). Next observe that \(f_*\) preserves the filtration (using results of Exposé VIII, which causes no difficulty, since 4.8 is not used to prove this fact). Indeed, it follows from 1.8 ii) that \(f\) is a complete-intersection morphism of virtual relative dimension zero (VIII 4.2). Since \(f\) is projective, the result follows from VIII 3.2. Applying \(f_*\) to the relation
